\documentclass[conference]{IEEEtran}
\IEEEoverridecommandlockouts
\usepackage{cite}
\usepackage{amsmath,amssymb}
\usepackage{graphicx}
\usepackage{booktabs}
\usepackage{url}
\usepackage{xcolor}
\usepackage[hidelinks]{hyperref}
\usepackage{listings}

\lstset{basicstyle=\ttfamily\footnotesize,breaklines=true,columns=fullflexible,
        frame=single,framesep=3pt,xleftmargin=2pt}

\begin{document}

\title{RippleKG: Evidence-Aware Incremental View\\Maintenance for LLM-Generated Knowledge Graphs}

\author{%
\IEEEauthorblockN{Group~XX\,---\,Author A, Author B, Author C, Author D}
\IEEEauthorblockA{Database Management Systems (DB114), National Taiwan University\\
\texttt{\{author A--D\}@example.ntu.edu.tw}}}

\maketitle

\begin{abstract}
Retrieval-augmented generation brought semantic search into the database;
this project brings \emph{incremental view maintenance} (IVM) into
LLM-generated knowledge graphs (KGs). When a source corpus is edited, an
LLM-derived KG must be updated. Re-extracting the whole graph is expensive, and
invalidating every object reachable from a changed sentence over-invalidates the
graph. RippleKG instead treats the KG as a materialized view over
sentence-level provenance and propagates a corpus edit only along the evidence
that actually changed. We contribute three mechanisms, all materialized in a
single ArangoDB instance: \textbf{M1}, a semantic evidence delta that classifies
each affected evidence record as added, removed, or unchanged using a canonical
triple key; \textbf{M2}, a cost-aware invalidation policy that chooses
\textsc{Skip}/\textsc{Patch}/\textsc{Rebuild} per affected object; and persisted
per-object \emph{freshness} state. On Re-DocRED, the incrementally maintained KG
is provably consistent with a from-scratch rebuild (0 mismatches over 1{,}880
relations and 1{,}033 entities), while marking $4.68\times$ fewer objects stale
than reachability baselines and costing up to $\sim$$3{,}500\times$ less than a
full corpus rebuild.
\end{abstract}

\begin{IEEEkeywords}
incremental view maintenance, knowledge graph, provenance, retrieval-augmented
generation, ArangoDB
\end{IEEEkeywords}

\section{Introduction}
Large language models (LLMs) are increasingly used to build knowledge graphs
from text corpora, and graph-based retrieval (GraphRAG) consumes those graphs
downstream. In any persistent deployment the source corpus keeps changing: a
sentence is corrected, a fact is replaced, a passage is rewritten. The derived
KG must follow. Two strategies dominate practice and both are unsatisfactory.
A \emph{full rebuild} re-extracts the entire KG after every edit and is the most
expensive option. \emph{Naive} or \emph{generic-traversal} invalidation marks
every object reachable from a changed sentence as stale; it is cheap to decide
but \emph{over-invalidates}, because graph reachability is far coarser than the
evidence that a specific edit actually touches. Existing incremental
GraphRAG systems (e.g., LightRAG~\cite{lightrag}, Microsoft
GraphRAG~\cite{graphrag}, iText2KG~\cite{itext2kg}) are append/merge oriented:
they grow the graph but rarely delete, invalidate, or expose freshness.

We observe that this is exactly the problem incremental view maintenance solves
in databases, recast for LLM-generated views. The KG is a \emph{materialized
view}; sentences are base tuples; provenance edges are the dependency graph; a
corpus edit is a base-table update that must propagate to derived objects. Our
one-line framing:

\begin{quote}\itshape
RAG brought semantic search into the database; RippleKG brings incremental view
maintenance into LLM-generated knowledge graphs.
\end{quote}

\noindent\textbf{Contributions.}
(1) We frame LLM-KG maintenance as IVM and identify why classic IVM does not
directly apply (Sec.~\ref{sec:why}). (2) We design and implement three
mechanisms---M1 semantic evidence delta, M2 cost-aware
\textsc{Skip}/\textsc{Patch}/\textsc{Rebuild} policy, and persisted freshness
(Sec.~\ref{sec:method})---entirely on ArangoDB, so that evidence deltas,
invalidation decisions, and freshness are durable database state, not transient
objects (Sec.~\ref{sec:system}). (3) We evaluate correctness, over-invalidation,
and cost against full-rebuild, generic-traversal, and naive baselines, and report
a retrieval/gate/end-to-end study with an independently audited label set
(Sec.~\ref{sec:eval}).

\section{Why Classic IVM Is Insufficient}
\label{sec:why}
Classic IVM assumes a view is a deterministic, exact function of its base data
and that refresh is cheap enough to run eagerly. LLM-generated views break all
three assumptions, as summarized in Table~\ref{tab:ivm}. Re-extraction may emit
textually different evidence for the same underlying claim (paraphrase), so
equality must be \emph{semantic}, not byte-exact. Dependencies are fuzzy and
mediated by provenance rather than exact column references. And refresh can be
expensive (an LLM call), so it must be selective and, ideally, deferrable. These
gaps motivate a semantic delta operator, a cost-aware policy, and visible
freshness state.

\begin{table}[t]
\caption{Classic IVM assumptions vs.\ LLM-generated views}
\label{tab:ivm}
\centering
\footnotesize
\begin{tabular}{p{0.46\columnwidth}p{0.46\columnwidth}}
\toprule
\textbf{Classic IVM assumes} & \textbf{LLM-generated view reality}\\
\midrule
View is deterministic & Re-extraction paraphrases identical meaning\\
Dependencies are exact & Dependencies are fuzzy / semantic\\
Refresh is cheap (eager) & Extraction/judging is costly; refresh must be selective\\
\bottomrule
\end{tabular}
\end{table}

\section{Related Work}
RippleKG sits at the intersection of four neighboring lines of work.
\emph{Incremental LLM-KG / GraphRAG} systems~\cite{graphrag,lightrag,itext2kg}
build and grow KGs from text. Microsoft GraphRAG~\cite{graphrag} rebuilds
community summaries, LightRAG~\cite{lightrag} performs dual-level incremental
insertion, and iText2KG~\cite{itext2kg} merges newly extracted triples into an
existing graph. All are append/merge oriented: they add or reconcile nodes but
rarely \emph{retract} evidence, invalidate stale objects, or expose per-object
freshness, so the corpus$\rightarrow$KG maintenance loop is left open.
\emph{Classic database IVM}~\cite{guptamumick,dbtoaster,diffdataflow} maintains
materialized views under deltas that are exact, deterministic, and cheap enough
for eager refresh---assumptions Sec.~\ref{sec:why} shows do not hold for an
LLM-generated view, motivating a semantic delta and a cost-aware policy rather
than algebraic delta rules. \emph{KG provenance and outdated-fact detection}
track which source supports which fact and flag staleness, but stop short of an
end-to-end, cost-aware refresh policy backed by persistent decision state.
Finally, work sharing the ``ripple'' name, notably RippleEdits~\cite{rippleedits},
studies ripple effects of editing \emph{model weights} in an LLM, an orthogonal
problem to maintaining a database-materialized KG. To our knowledge, the
combination of a semantic evidence delta, a cost-aware
\textsc{Skip}/\textsc{Patch}/\textsc{Rebuild} policy, and persisted per-object
freshness for a database-materialized LLM-KG is unoccupied.

\section{Method}
\label{sec:method}

\subsection{Pipeline and the synchronous spine}
A corpus edit is supplied as an \texttt{EditOp}$=(\textit{doc\_id},
\textit{sent\_idx}, \textit{new\_text}, \textit{intended\_triples})$. We adopt
\emph{authored intended triples} (the set of triples the edited sentence should
support afterwards) as the primary evidence source, so that extraction quality
is not the object of study; a lightweight LLM extractor can produce the same
interface. Processing follows six stages:

\begin{lstlisting}
corpus edit
 -> changed sentence (text_hash)
 -> affected evidence (1-hop provenance)
 -> M1 semantic evidence delta
 -> affected entity / relation
 -> M2 SKIP / PATCH / REBUILD
 -> persisted KG / provenance / freshness
\end{lstlisting}

Stages 1--5 form a \emph{synchronous spine} executed inside one ArangoDB stream
transaction: the sentence text, provenance edges, and the existence of entities
and relations are always made consistent atomically. Only the derived-aggregate
repair (recomputing evidence counts, dropping empty relations, flipping freshness
back) may be deferred (Sec.~\ref{sec:fresh}).

\subsection{M1: semantic evidence delta}
Given a changed sentence, M1 collects its old active evidence by a 1-hop outbound
traversal over the provenance edges and compares it to the new
\texttt{intended\_triples}. Relation evidence is keyed by the canonical triple
$(\textit{head}, \textit{rel\_type}, \textit{tail})$ after name normalization;
mention evidence is keyed by the normalized entity name:
\[
\textit{new}\setminus\textit{old}\!\rightarrow\!\textsc{added},\;
\textit{old}\setminus\textit{new}\!\rightarrow\!\textsc{removed},\;
\textit{old}\cap\textit{new}\!\rightarrow\!\textsc{unchanged}.
\]
The crucial case is \textsc{unchanged}: when only the surface wording changes but
the canonical triples are identical (a paraphrase), M1 reports no change---which
exact, hash-based IVM cannot do. M1 simultaneously updates provenance edges
(insert for \textsc{added}, mark \texttt{status=removed} for \textsc{removed},
leave \textsc{unchanged} untouched) and persists each change to the
\texttt{evidence\_deltas} collection.

\subsection{M2: cost-aware invalidation policy}
For each affected object M2 consults the decision table in
Table~\ref{tab:policy}. The policy is the cost-awareness: \textsc{Skip} when
nothing changed, \textsc{Patch} when an aggregate can be repaired in place, and
\textsc{Rebuild} only when the last supporting evidence disappears. Nominal costs
are $\textsc{Skip}=0$, $\textsc{Patch}=1$ (a database write), and
$\textsc{Rebuild}=5$ (modeling re-resolution/re-extraction as $5\times$ a patch).
Each decision is written to \texttt{refresh\_decisions} with
\texttt{status=pending}.

\begin{table}[t]
\caption{M2 decision table}
\label{tab:policy}
\centering
\footnotesize
\begin{tabular}{p{0.58\columnwidth}ll}
\toprule
\textbf{Evidence delta on object} & \textbf{Decision} & \textbf{Cost}\\
\midrule
Triple unchanged (paraphrase) & \textsc{Skip} & 0\\
Evidence added & \textsc{Patch} & 1\\
Evidence removed, others remain active & \textsc{Patch} & 1\\
Last active evidence removed & \textsc{Rebuild} & 5\\
\bottomrule
\end{tabular}
\end{table}

\subsection{Freshness and refresh execution}
\label{sec:fresh}
Every entity and relation carries \texttt{freshness\_status}
(\texttt{fresh}/\texttt{stale}) and \texttt{last\_changed\_step}; history is
reconstructable from the delta/decision logs (no bitemporal versioning).
\texttt{apply\_refreshes} processes pending decisions in two modes.
\emph{Immediate} (default) refreshes within the step, so all objects return to
\texttt{fresh}; \emph{deferred} leaves chosen decisions pending until a later
maintenance tick, making staleness visible across steps. Deferral is safe by
three invariants: (i) the next edit's M1 does not depend on aggregates;
(ii) \texttt{apply\_refreshes} is idempotent (it always recomputes from current
evidence); and (iii) stale objects are explicitly flagged, i.e., \emph{known}
not-yet-refreshed rather than silently wrong. Queries may request
\texttt{fresh\_only} to exclude stale objects.

\section{System Design and Schema}
\label{sec:system}
RippleKG is implemented on a single ArangoDB database; a FastAPI service imports
the package as a library so the demo, scripts, and notebook all call the same
\texttt{run\_edit} path and never drift from the real pipeline. Because we do not
modify storage internals, the contribution is framed as \emph{logical database
design}: a provenance schema, evidenced records, AQL provenance traversal,
transactional invalidation, and visible freshness. Table~\ref{tab:map} maps each
pipeline element to its DBMS concept.

\begin{table}[t]
\caption{Pipeline$\rightarrow$DBMS concept mapping}
\label{tab:map}
\centering
\footnotesize
\begin{tabular}{p{0.46\columnwidth}p{0.46\columnwidth}}
\toprule
\textbf{RippleKG element} & \textbf{DBMS concept}\\
\midrule
corpus update & base-table update\\
changed sentence / diff & change capture\\
provenance edges & forward dependency tracking\\
evidence delta (M1) & view-level change detection\\
affected entity/relation & derived-table propagation\\
\textsc{Skip}/\textsc{Patch}/\textsc{Rebuild} (M2) & cost-aware view maintenance\\
\texttt{freshness\_status} & materialized-view freshness\\
\bottomrule
\end{tabular}
\end{table}

The store uses eight collections: four document collections
(\texttt{documents}, \texttt{sentences}, \texttt{entities}, \texttt{relations}),
two provenance edge collections along which the ripple travels
(\texttt{mentions}: sentence$\rightarrow$entity;
\texttt{sentence\_supports\_relation}: sentence$\rightarrow$relation), and two
maintenance collections (\texttt{evidence\_deltas} for M1 output and
\texttt{refresh\_decisions} for M2 output). Because provenance is a real graph,
the affected set is exactly a 1-hop outbound traversal from the changed
sentence---no multi-hop reachability is needed---which is the structural reason
RippleKG avoids over-invalidation.

\section{Evaluation}
\label{sec:eval}

\subsection{Setup}
We build the initial $T_0$ graph directly from Re-DocRED~\cite{redocred}
annotations (a deterministic corpus$\rightarrow$KG pair derived from
DocRED~\cite{docred}) without running an LLM: \texttt{sents} become sentences,
\texttt{vertexSet} becomes entities and mention edges, \texttt{labels} become
relations, and \texttt{labels.evidence} becomes supporting edges. Synthetic edits
carry authored \texttt{intended\_triples}. The store is ArangoDB~3.12; embeddings
(optional) use \texttt{all-MiniLM-L6-v2}~\cite{sbert}.

\subsection{Case studies}
Two representative single-sentence edits exercise the policy paths
(Table~\ref{tab:case}). doc0's opening sentence states that Schneider was born
in Media\c{s}, Transylvania, and is a \emph{German} skeleton racer. In
\emph{Case~A} we reorder it into ``Born on 13 March 1963 in
Media\c{s}\dots{}, Schneider is a German skeleton racer'': the canonical
triples are identical, so M1 returns only \textsc{Unchanged} deltas, all 12
affected objects \textsc{Skip}, and the incremental cost is zero against a
full-document rebuild of nominal cost~60. In \emph{Case~B} we rewrite the
birthplace and nationality to \emph{Krak\'ow, Poland} and \emph{Polish}; M1
reports the old place-of-birth and citizenship triples \textsc{removed} and the
new ones \textsc{added} (date of birth \textsc{Unchanged}), and M2
\textsc{Rebuild}s the relations whose last supporting evidence vanished while
\textsc{Patch}ing those that retain evidence---touching only the affected
objects at $\sim$$49\%$ of a full-document rebuild. Crucially, Case~A is exactly
where reachability baselines over-invalidate: they would mark all 12 mentioned
objects stale, whereas evidence-aware comparison correctly skips them.

\begin{table}[t]
\caption{Two representative edits on doc0 (A: paraphrase, B: factual change)}
\label{tab:case}
\centering
\footnotesize
\setlength{\tabcolsep}{4pt}
\begin{tabular}{lcc}
\toprule
& \textbf{Case A} & \textbf{Case B}\\
\midrule
M1 (unchanged) & 12 & 3\\
M1 (added / removed) & 0 / 0 & 4 / 9\\
M2 \textsc{Skip}/\textsc{Patch}/\textsc{Rebuild} & 12 / 0 / 0 & 3 / 6 / 7\\
incremental cost & 0 & 41\\
full-doc rebuild cost & 60 & 80\\
reduction & 100\% & 48.75\%\\
\bottomrule
\end{tabular}
\end{table}

\subsection{Baselines: correctness, over-invalidation, cost}
We compare against three baselines (all optional sanity references): \textbf{B0}
full rebuild (recompute every object's aggregate from active evidence), \textbf{B1}
generic-traversal invalidation (mark all reachable objects stale), and \textbf{B2}
naive invalidation (mark every mentioned object stale). Running 48 mixed edits
over 50 documents (Table~\ref{tab:bench}):

\emph{Correctness (B0).} After all incremental edits, the materialized
\texttt{evidence\_count}/\texttt{status} of every object equals a from-scratch
recomputation over active provenance edges: \textbf{0 mismatches} across 1{,}880
relations and 1{,}033 entities. This is the classic IVM consistency invariant;
incremental \textsc{Skip}/\textsc{Patch}/\textsc{Rebuild} did not sacrifice
correctness.

\emph{Over-invalidation (B1/B2).} RippleKG marks 90 objects stale; both
reachability baselines mark 421, i.e.\ $4.68\times$ more. The gap comes from
paraphrase and non-evidence edits, which RippleKG \textsc{Skip}s but the
baselines still invalidate.

\emph{Cost (three granularities).} Nominal cost in \textsc{Rebuild} units shows
the value of fine granularity: per-sentence (ours) $\ll$ per-document
(GraphRAG-style) $\ll$ whole-corpus rebuild.

\begin{table}[t]
\caption{Baseline comparison (50 docs, 48 mixed edits)}
\label{tab:bench}
\centering
\footnotesize
\begin{tabular}{lrr}
\toprule
\textbf{Metric} & \textbf{Value} & \textbf{vs.\ ours}\\
\midrule
Decisions \textsc{Skip}/\textsc{Patch}/\textsc{Rebuild} & 355/63/27 & ---\\
\midrule
Objects marked stale (ours) & 90 & $1\times$\\
\quad B1 generic-traversal & 421 & $4.68\times$\\
\quad B2 naive & 421 & $4.68\times$\\
\midrule
Cost: ours (sentence) & 198 & $1\times$\\
\quad document rebuild (GraphRAG-style) & 18{,}730 & $94.6\times$\\
\quad whole-KG rebuild (B0) & 693{,}750 & $3503.8\times$\\
\midrule
B0 correctness & \multicolumn{2}{r}{consistent, 0 mismatches}\\
\bottomrule
\end{tabular}
\end{table}

\subsection{Retrieval, relevance gate, and end-to-end}
A hybrid front end can select \emph{which} sentences to edit before M1/M2 run.
On a 50-fact evidence-retrieval benchmark, Transformer retrieval reaches
Hit@1 $=86\%$, Hit@3 $=100\%$, MRR $=0.93$; graph provenance gives a high-recall
(R$=100\%$) but lower-precision (P$=58.8\%$) affected set. An LLM relevance gate
raises precision to $76.9\%$ at R$=100\%$ (F1 $=87\%$). End-to-end, after adding a
deterministic replacement-aware verifier in front of M1, EditOp content accuracy
is $95\%$ and M1 added/removed accuracy and end-to-end M2 accuracy both reach
$100\%$ (Table~\ref{tab:e2e}); M2 policy-rule accuracy is $100\%$ throughout,
confirming that residual end-to-end error originated in the intended-triple
verifier, not the policy.

\begin{table}[t]
\caption{End-to-end accuracy (20 should-edit operations)}
\label{tab:e2e}
\centering
\footnotesize
\begin{tabular}{lcc}
\toprule
\textbf{Stage} & \textbf{before verifier} & \textbf{after}\\
\midrule
EditOp content accuracy & 80\% & \textbf{95\%}\\
M1 added/removed accuracy & 80\% & \textbf{100\%}\\
End-to-end M2 accuracy & 80\% & \textbf{100\%}\\
M2 policy-rule accuracy & 100\% & 100\%\\
\bottomrule
\end{tabular}
\end{table}

\subsection{Label audit}
The should-edit gold set (30 replacement cases, 150 candidate sentences) was
independently re-checked against its annotation rule: of 150 labels,
\textbf{0 were overturned}; all 20 \textsc{True} labels genuinely assert the old
fact, and 12 borderline \textsc{False} labels (the value appears in the sentence
but is not asserted, e.g.\ ``the \emph{Canadian} province of Alberta'') were
confirmed correct. The audit notably validates that the same word in one sentence
is labeled \textsc{True} for the fact it asserts and \textsc{False} for a fact it
merely mentions.

\section{Limitations}
Evidence is supplied as authored triples (option A); an LLM extractor (option B)
shares the same interface but is left to future work. Costs are nominal units,
not measured dollar/token cost. The reviewed benchmark is small (5--50 documents,
a single \texttt{replace\_old\_fact} update type) and some synthetic replacements
are semantically unnatural; the labels handle these correctly but the inputs are
noisy. Query-time lazy refresh, \texttt{max\_staleness}, and entity-resolution
ripple are designed but not implemented.

\section{Conclusion}
RippleKG reframes the maintenance of an LLM-generated knowledge graph as
incremental view maintenance and realizes it entirely as durable ArangoDB state:
a semantic evidence delta (M1), a cost-aware
\textsc{Skip}/\textsc{Patch}/\textsc{Rebuild} policy (M2), and persisted per-object
freshness. A corpus edit ripples only along the evidence that actually changed,
yielding a KG provably consistent with a full rebuild while marking far fewer
objects stale and costing orders of magnitude less. RAG brought semantic search
into the database; RippleKG brings incremental view maintenance into
LLM-generated knowledge graphs.

\section{Member Contributions}
The division of labor follows the ownership structure fixed in the project
proposal. The percentages reflect each member's overall effort and sum to
100\%.
% TODO(team): replace <name> with real names/IDs and adjust the percentages
% so they reflect each member's ACTUAL contribution (must still total 100).
\begin{table}[h]
\caption{Per-member contribution}
\label{tab:contrib}
\centering
\footnotesize
\setlength{\tabcolsep}{4pt}
\begin{tabular}{p{0.18\columnwidth}p{0.56\columnwidth}c}
\toprule
\textbf{Member (role)} & \textbf{Primary contributions} & \textbf{\%}\\
\midrule
$\langle$name$\rangle$ (A) & Re-DocRED loader and $T_0$ build, synthetic edit set, B0/B1/B2 baselines & 25\\
$\langle$name$\rangle$ (B) & ArangoDB schema (8 collections), indexes and provenance persistence, AQL 1-hop traversal and DB inspection & 25\\
$\langle$name$\rangle$ (C) & M1 semantic evidence delta, M2 cost-aware policy, refresh execution & 25\\
$\langle$name$\rangle$ (D) & Evaluation harness and metrics, label audit, report, slides, demo and video & 25\\
\midrule
\textbf{Total} & & \textbf{100}\\
\bottomrule
\end{tabular}
\end{table}

\begin{thebibliography}{10}
\footnotesize
\bibitem{docred} Y.~Yao \emph{et al.}, ``DocRED: A Large-Scale Document-Level
Relation Extraction Dataset,'' in \emph{ACL}, 2019.
\bibitem{redocred} Q.~Tan \emph{et al.}, ``Revisiting DocRED---Addressing the
False Negative Problem in Relation Extraction,'' in \emph{EMNLP}, 2022.
\bibitem{graphrag} D.~Edge \emph{et al.}, ``From Local to Global: A Graph RAG
Approach to Query-Focused Summarization,'' \emph{arXiv:2404.16130}, 2024.
\bibitem{lightrag} Z.~Guo \emph{et al.}, ``LightRAG: Simple and Fast
Retrieval-Augmented Generation,'' \emph{arXiv:2410.05779}, 2024.
\bibitem{itext2kg} Y.~Lairgi \emph{et al.}, ``iText2KG: Incremental Knowledge
Graphs Construction Using Large Language Models,'' \emph{arXiv:2409.03284}, 2024.
\bibitem{rippleedits} R.~Cohen \emph{et al.}, ``Evaluating the Ripple Effects of
Knowledge Editing in Language Models,'' \emph{TACL}, 2024.
\bibitem{guptamumick} A.~Gupta and I.~S.~Mumick, ``Maintenance of Materialized
Views: Problems, Techniques, and Applications,'' \emph{IEEE Data Eng.\ Bull.}, 1995.
\bibitem{dbtoaster} C.~Koch \emph{et al.}, ``DBToaster: Higher-order Delta
Processing for Dynamic, Frequently Fresh Views,'' in \emph{VLDB}, 2014.
\bibitem{diffdataflow} F.~McSherry \emph{et al.}, ``Differential Dataflow,'' in
\emph{CIDR}, 2013.
\bibitem{sbert} N.~Reimers and I.~Gurevych, ``Sentence-BERT: Sentence Embeddings
using Siamese BERT-Networks,'' in \emph{EMNLP}, 2019.
\end{thebibliography}

\end{document}
